\documentclass[11pt]{article}
\usepackage[margin=1in]{geometry}
\usepackage{amsmath,amssymb,amsthm}
\usepackage{xurl}
\usepackage{hyperref}
\sloppy
\let\oldus\_
\renewcommand{\_}{\oldus\allowbreak}

\newtheorem{theorem}{Theorem}
\newtheorem{lemma}[theorem]{Lemma}
\theoremstyle{definition}
\newtheorem{definition}[theorem]{Definition}
\theoremstyle{remark}
\newtheorem{remark}[theorem]{Remark}

\newcommand{\Der}{\operatorname{Der}}
\newcommand{\Inn}{\operatorname{Inn}}
\newcommand{\ad}{\operatorname{ad}}
\newcommand{\HH}{\operatorname{HH}}
\newcommand{\lean}[1]{\texttt{#1}}

\title{Conjecture 00000004044 is false:\\ $\HH^1$ of the Kronecker quiver $K_3$ is not solvable}
\author{Jackmeson1}
\date{}

\begin{document}
\maketitle

\section{The conjecture and its reading}

Conjecture 00000004044 (English text):
\begin{quote}
Definition: HH\^{}1 is the Lie algebra of outer derivations, the first Hochschild cohomology.
Conjecture: The HH\^{}1 of a quiver algebra is solvable if and only if its oriented cycles are all
broken after removing at most one vertex; the decision table is complete for $n \le 6$ and the
formula holds for all $n$.
\end{quote}
The Chinese title also says ``solvability criterion for $\HH^1$'' (\emph{ke jie xing}), but the
Chinese body uses the word \emph{ke yue} (``reducible'') in the biconditional; see
Remark~\ref{rem:reading}.

\paragraph{Reading.} For a quiver $Q$ and a field $k$, the claim is
\[
  \HH^1(kQ/I) \text{ is solvable} \iff \text{some set of at most one vertex meets every oriented cycle of } Q .
\]
We refute the direction ``$\Leftarrow$'' with a single quiver algebra, over every field $k$, with
$n = 2$ vertices (so it lies inside the range $n \le 6$ of the claimed decision table).

\paragraph{Counterexample.} The generalized Kronecker quiver $K_3$ has vertices $1,2$ and three
parallel arrows $\alpha_0,\alpha_1,\alpha_2 : 1 \to 2$. It has no oriented cycle, so the cycle
condition holds with zero vertices removed. Its path algebra $kK_3$ (ideal $I=0$) is a quiver
algebra in every usual sense: it is finite dimensional, and since $K_3$ has no path of length $2$,
$I = 0 = R_Q^2$ is an admissible ideal. We prove $\HH^1(kK_3)$ is not solvable.
(It is known that $\HH^1(kK_3) \cong \mathfrak{pgl}_3(k)$; see the quote from [LRW] below. We do
not use this isomorphism; we prove non-solvability directly.)

\section{Definitions (as formalized)}

All definitions are in the Lean file \lean{Conjecture4044/Basic.lean} (namespace \lean{C4044}).

\begin{definition}[Derivations, inner derivations, $\HH^1$]
Let $A$ be an associative unital algebra over a commutative ring $k$.
\lean{Der k A} is the set of $k$-linear maps $D : A \to A$ with $D(xy) = D(x)y + xD(y)$, a Lie
subalgebra of $\mathrm{End}_k(A)$ under the commutator $[D,E] = D\circ E - E \circ D$.
For $a \in A$, $\ad a : x \mapsto ax - xa$ is a derivation (\lean{ad\_mem\_Der}).
\lean{Inn k A} $= \{\ad a : a \in A\}$ is a Lie ideal of \lean{Der k A} (the ideal property is
$[D, \ad a] = \ad(D a)$). Finally \lean{HH1 k A} $:= \Der(A)/\Inn(A)$ with Mathlib's quotient Lie
algebra structure.
Mathlib's \lean{Derivation} requires a commutative algebra, so these are defined directly.
Solvability is Mathlib's \lean{LieAlgebra.IsSolvable}: some term of the derived series
$L^{(0)} = L$, $L^{(m+1)} = [L^{(m)}, L^{(m)}]$ is $0$.
\end{definition}

This is the standard description of $\HH^1(A) = \HH^1(A,A)$: [LRW, Section 2] says (math symbols
stripped by the retrieval tool) ``the elements in ( ) can be interpreted as the -derivations of
into and the elements in ( ) can be interpreted as the inner -derivations of into . It is
well-known that ( ) is a Lie algebra under the Lie bracket and ( ) is a Lie ideal of ( ), so that
( ) has a Lie algebra structure.''

\begin{definition}[The quiver and the cycle condition]
\lean{V} $=\{1,2\}$ with Mathlib \lean{Quiver} structure \lean{Hom 1 2 = Fin 3} and all other hom
types empty. A path $p$ \lean{visits} $v$ if $v$ is its source or the target of one of its arrows.
\lean{CyclesBrokenByAtMostOneVertex W} means: there is a set $S$ of vertices with at most one
element (\lean{S.Subsingleton}) such that every closed path of positive length visits some
vertex of $S$; equivalently, deleting $S$ leaves no oriented cycle.
\end{definition}

\begin{definition}[The path algebra $kK_3$]
\lean{P} is the set of paths of $K_3$: trivial paths $e_1,e_2$ and the arrows $\alpha_i$.
\lean{paths\_bijective} proves that \lean{P.toPath} is a bijection onto all paths
$\Sigma\, a\, b,\ \lean{Quiver.Path}\ a\ b$ of the Mathlib quiver.
\lean{PathAlg k} $= k^{P}$ is the vector space with basis $P$, with the product given on basis
elements by concatenation (\lean{basis\_mul\_basis}): $e_1e_1 = e_1$, $e_2e_2=e_2$,
$e_1\alpha_i = \alpha_i = \alpha_i e_2$, and every other product of two paths is $0$ (the paths are
not composable). Paths compose left to right, as in Mathlib's \lean{Quiver.Path.comp}.
This is the path algebra as defined in [Wiki]: ``multiplication given by concatenation of paths.
If two paths cannot be concatenated because the end vertex of the first is not equal to the
starting vertex of the second, their product is defined to be zero. This defines an associative
algebra over''. The ring and $k$-algebra axioms are proved in Lean.
\end{definition}

\section{Proof}

\begin{lemma}[\lean{K3\_no\_oriented\_cycle}, \lean{K3\_cyclesBroken}]
Every closed path in $K_3$ has length $0$. Hence $K_3$ satisfies the cycle condition with
$S = \emptyset$.
\end{lemma}
\begin{proof}
By induction on paths (\lean{path\_shape}), every path has length $0$, or goes from $1$ to $2$
and has length $1$; there is no arrow out of vertex $2$. A closed path of length $1$ would need
$1 = 2$.
\end{proof}

For $i \ne j$ in $\{0,1,2\}$ let $D_{ij}$ (\lean{Dlin i j}, \lean{D i j}) be the linear map with
$D_{ij}(\alpha_j) = \alpha_i$ and $D_{ij}$ zero on $e_1, e_2$ and on the other arrows. In
coordinates, $D_{ij}(x)$ has $\alpha_l$-coordinate $x_{\alpha_j}$ if $l = i$ and $0$ otherwise, and
zero $e$-coordinates.

\begin{lemma}[\lean{Dlin\_mem}]
Each $D_{ij}$ is a derivation of $kK_3$.
\end{lemma}
\begin{proof}
The $\alpha_l$-coordinate of $xy$ is $x_{e_1}y_{\alpha_l} + x_{\alpha_l}y_{e_2}$ and the
$e$-coordinates of $xy$ involve only $e$-coordinates. Both sides of $D(xy) = D(x)y + xD(y)$
then have $\alpha_l$-coordinate $[l=i](x_{e_1}y_{\alpha_j} + x_{\alpha_j}y_{e_2})$ and zero
$e$-coordinates.
\end{proof}

\begin{lemma}[\lean{D\_bracket}]
For distinct $i, j, m$: $[D_{im}, D_{mj}] = D_{ij}$ in $\Der(kK_3)$.
\end{lemma}
\begin{proof}
$D_{im}D_{mj}$ sends $\alpha_j \mapsto \alpha_i$ and kills everything else; $D_{mj}D_{im}$ is $0$
because $D_{im}$ has image in $k\alpha_i$ and $D_{mj}(\alpha_i) = 0$ for $i \ne j$.
\end{proof}

\begin{lemma}[\lean{D\_not\_inner}]
For $i \ne j$, $D_{ij} \notin \Inn(kK_3)$.
\end{lemma}
\begin{proof}
For any $a$, the $\alpha_i$-coordinate of $(\ad a)(\alpha_j) = a\alpha_j - \alpha_j a$ is
$a_{e_1}\cdot 0 + a_{\alpha_i}\cdot 0 - 0\cdot a_{\alpha_i} - 0 \cdot a_{e_2} = 0$, while that of
$D_{ij}(\alpha_j) = \alpha_i$ is $1$.
\end{proof}

\begin{theorem}[\lean{HH1\_not\_solvable}]
For every field $k$, the Lie algebra $\HH^1(kK_3) = \Der(kK_3)/\Inn(kK_3)$ is not solvable.
\end{theorem}
\begin{proof}
Let $d_{ij}$ (\lean{d i j}) be the class of $D_{ij}$. The quotient map is a Lie homomorphism, so
$[d_{im}, d_{mj}] = d_{ij}$ for distinct $i,j,m$ (\lean{d\_bracket}). By induction on $n$, every
$d_{ij}$ with $i\ne j$ lies in the $n$-th derived term (\lean{d\_mem\_derivedSeries}): for the
step pick $m \notin \{i,j\}$, then $d_{ij} = [d_{im}, d_{mj}]$ with both factors in the $n$-th term.
If the $n$-th term were $0$, then $d_{01} = 0$, i.e.\ $D_{01} \in \Inn$, contradicting
\lean{D\_not\_inner}. No characteristic assumption is used.
\end{proof}

\begin{theorem}[\lean{conjecture\_00000004044\_false}]
\lean{CyclesBrokenByAtMostOneVertex V $\wedge$ $\neg$ LieAlgebra.IsSolvable (HH1 k (PathAlg k))}
for every field \lean{k}. Consequently (\lean{conjecture\_00000004044\_iff\_fails})
\lean{$\neg$ (LieAlgebra.IsSolvable (HH1 k (PathAlg k)) $\leftrightarrow$
CyclesBrokenByAtMostOneVertex V)}.
\end{theorem}

\begin{remark}[Scope and conventions]\label{rem:reading}
(1) What is refuted: the ``solvable'' reading, for quiver algebras meaning either the path algebra
$kQ$ or a bound quiver algebra $kQ/I$ with $I$ admissible (both include $kK_3$), over any fixed
field. The ``only if'' direction is not addressed (one direction failing refutes the biconditional).
(2) The Chinese body says \emph{ke yue} (``reducible'') where the English body and the Chinese
title say solvable. We read solvable as the intended statement: two of the three texts say so, and
the conjecture's own parenthetical name is ``solvability criterion''. ``Reducible Lie algebra''
is not a standard notion; under the reading ``reducible = has a nonzero proper ideal'' the
statement is \emph{not} refuted here ($\HH^1(kK_3) \cong \mathfrak{sl}_3$ is simple in
characteristic $0$, but proving that needs the full $\HH^1$ computation and is not formalized).
(3) Conventions: paths compose left to right; $\HH^1$ uses the commutator bracket on
$\mathrm{End}_k(A)$ (Mathlib's \lean{LieRing.ofAssociativeRing}, enabled as a local instance as in
Mathlib's own Lie theory files); ``removing at most one vertex'' allows removing none.
\end{remark}

\section{Lean audit}

File \lean{lean/Conjecture4044/Basic.lean} (363 lines, only \lean{import Mathlib}), Lean
\lean{v4.33.1}, Mathlib \lean{v4.33.1}. Main theorem \lean{C4044.conjecture\_00000004044\_false}.
\lean{\#print axioms} for \lean{conjecture\_00000004044\_false},
\lean{conjecture\_00000004044\_iff\_fails} and \lean{paths\_bijective} gives
\lean{[propext, Classical.choice, Quot.sound]}; no \lean{sorry}, no \lean{native\_decide}, no
added axioms.

\section*{References}
\begin{itemize}
\item[{[LRW]}] Y. Liu, L. Rubio y Degrassi, C. Wen, \emph{On the first Hochschild cohomology of
finite dimensional quiver algebras under gluing idempotents}, \url{https://arxiv.org/abs/2211.05435v3}.
On the Kronecker quiver (math stripped by retrieval): ``if is the -Kronecker quiver (the quiver
with parallel arrows), with the convention that -Kronecker quiver is the Dnykin quiver where is
defined as . In other words, is the quotient by its center.''
\item[{[Wiki]}] Quiver (mathematics), \url{https://en.wikipedia.org/wiki/Quiver_(mathematics)}.
\end{itemize}

\end{document}
